\section{Smoothness of the level sets}\label{sec:smooth}

Throughout this section $\Om\subset\R^2$ is a bounded open set that is
\emph{regular for the Dirichlet problem}, and $u$ denotes its torsion
function: the unique
\[
  u\in H^1_0(\Om)\cap C(\overline\Om),\qquad -\Delta u=1\ \text{in }\Om,
  \qquad u=0\ \text{on }\partial\Om .
\]
We collect from the Preliminaries (Section~\ref{sec:prelim}) the four facts
we are allowed to use without reproof:
\begin{enumerate}
\item[(P1)] $u$ is real-analytic in $\Om$ (interior analyticity of solutions
  of $-\Delta u=1$; see \cite[Ch.~6, \S6.1 and the analyticity
  theorem]{Morrey1966} or \cite[Theorem~6.6.1]{Morrey1966}).
\item[(P2)] Interior gradient and Hessian estimates: for every $K\Subset\Om$
  there is $C(K)$ with $\|u\|_{C^2(K)}\le C(K)$
  (\cite[Theorems~6.2, 6.17 and Corollary~6.3]{GT1983}).
\item[(P3)] Sard's theorem: the set of critical values of a $C^k$ map
  $\R^2\to\R$ with $k\ge2$ has Lebesgue measure zero
  (\cite[\S2--3]{Milnor1997}; \cite{Sard1942}).
\item[(P4)] The coarea formula and the structure of sets of finite perimeter
  (Section~\ref{sec:prelim}; \cite[\S3.4, \S5]{EvansGariepy2015}).
\end{enumerate}

\medskip
\noindent\textbf{Positivity of $u$.}
Since $-\Delta u=1>0$ in $\Om$, $u$ is superharmonic and not constant; as
$u=0$ on $\partial\Om$ and $u\not\equiv0$, the strong maximum principle
(equivalently Hopf's lemma) forbids an interior minimum equal to the boundary
value unless $u$ is constant, hence $u>0$ throughout the connected
components of $\Om$ meeting the interior; concretely, if $u(x_0)=0=\min u$ at
some $x_0\in\Om$ then the strong maximum principle for $-\Delta u\ge0$ would
force $u\equiv0$ near $x_0$, contradicting $-\Delta u=1$. Therefore $u>0$ in
$\Om$. Set
\[
  m:=\max_{\overline\Om}u=\max_{\Om}u>0 ,
\]
the maximum being attained by continuity on the compact set $\overline\Om$
and being positive since $u>0$ somewhere. For $v\in(0,m)$ write
\[
  E_v:=\{x\in\Om:u(x)>v\} .
\]

\subsection*{Compact containment of superlevel sets}

\begin{lemma}[Compact containment]\label{lem:compact-containment}
For every $v\in(0,m)$ the superlevel set
$\overline{E_v}=\{x\in\overline\Om:u(x)\ge v\}=\{u\ge v\}$ is a nonempty
compact subset of the open set $\Om$, and
\[
  \operatorname{dist}\bigl(\{u\ge v\},\partial\Om\bigr)\ \ge\ c(v)>0 .
\]
\end{lemma}

\begin{proof}
Since $u\in C(\overline\Om)$ and $\overline\Om$ is compact, the function
$u\colon\overline\Om\to\R$ is continuous on a compact set.

\emph{Nonempty.} By definition of $m$ there is $x_*\in\Om$ with $u(x_*)=m>v$,
so $x_*\in\{u\ge v\}$.

\emph{Closed and bounded.} The set $\{u\ge v\}=u^{-1}([v,\infty))$ is the
preimage of a closed set under the continuous map $u$ on $\overline\Om$, hence
relatively closed in $\overline\Om$; as $\overline\Om$ is itself closed in
$\R^2$, $\{u\ge v\}$ is closed in $\R^2$. It is contained in the bounded set
$\overline\Om$, hence bounded. A closed and bounded subset of $\R^2$ is
compact (Heine--Borel).

\emph{Disjoint from $\partial\Om$.} On $\partial\Om$ we have $u=0<v$, so no
boundary point lies in $\{u\ge v\}$. Thus $\{u\ge v\}\subset\Om$.

\emph{Positive distance.} The two sets $\{u\ge v\}$ and $\partial\Om$ are
disjoint and both compact (the first by the above, the second because
$\partial\Om$ is closed and bounded). The distance function
$x\mapsto\operatorname{dist}(x,\partial\Om)$ is continuous and strictly
positive on the compact set $\{u\ge v\}$ (strictly positive because that set
is disjoint from the closed set $\partial\Om$), so it attains a positive
minimum
\[
  c(v):=\min_{x\in\{u\ge v\}}\operatorname{dist}(x,\partial\Om)>0 .
\]
This is the asserted lower bound. Finally
$\overline{E_v}=\{u\ge v\}$: clearly $E_v\subset\{u\ge v\}$ and the latter is
closed, so $\overline{E_v}\subset\{u\ge v\}$; conversely if $u(x)\ge v$ then,
$u$ being continuous and nonconstant near $x$ (analytic, $-\Delta u=1$), every
neighbourhood of $x$ meets $\{u>v\}$, whence $x\in\overline{E_v}$. (We will
re-derive $\partial E_v=\{u=v\}$ more precisely in
Lemma~\ref{lem:smooth-level} for regular $v$.)
\end{proof}

\begin{remark}\label{rem:uniform-c2}
Combining Lemma~\ref{lem:compact-containment} with (P2): fixing
$v\in(0,m)$, the compact set $K_v:=\{u\ge v/2\}\Subset\Om$ contains an open
neighbourhood of $\{u\ge v\}$, and on $K_v$ all estimates (P1)--(P2) apply.
In particular $u\in C^\infty(K_v)$, indeed real-analytic, with $\nabla u$ and
$D^2u$ bounded there. This is what lets us apply the divergence theorem and
the implicit function theorem in a fixed interior region.
\end{remark}

\subsection*{Regular level sets are analytic curves}

Recall that $v\in(0,m)$ is a \emph{regular value} of $u$ if $\nabla u(x)\ne0$
for every $x\in\{u=v\}$; otherwise $v$ is a \emph{critical value}.

\begin{lemma}[Smooth regular level sets]\label{lem:smooth-level}
Let $v\in(0,m)$ be a regular value of $u$. Then:
\begin{enumerate}
\item[(a)] $\{u=v\}$ is a nonempty compact subset of $\Om$ on which
  $|\nabla u|>0$ everywhere;
\item[(b)] $\{u=v\}$ is a real-analytic embedded $1$-manifold without
  boundary, equal to a finite disjoint union of analytic Jordan curves
  $\Gamma_1,\dots,\Gamma_N$;
\item[(c)] $E_v=\{u>v\}$ is a bounded open subset of $\Om$ with topological
  boundary $\partial E_v=\{u=v\}$;
\item[(d)] $E_v$ has finite perimeter,
  $P(E_v)=\HH^1(\{u=v\})=\sum_{i=1}^N\HH^1(\Gamma_i)$, and its measure-theoretic
  outward unit normal is
  \[
    \nu(x)=-\frac{\nabla u(x)}{|\nabla u(x)|}\qquad\text{for every }x\in\{u=v\}.
  \]
\end{enumerate}
\end{lemma}

\begin{proof}
\emph{Step 0: $\{u=v\}$ is compact and nonempty.}
By Lemma~\ref{lem:compact-containment}, $\{u\ge v\}\Subset\Om$ is compact, and
$\{u=v\}=\{u\ge v\}\cap\{u\le v\}$ is closed (intersection of two closed sets),
hence compact as a closed subset of a compact set. It is nonempty: $u$ attains
both the value $m>v$ (at the interior maximum) and, on $\partial\Om$, the value
$0<v$; by the intermediate value theorem along any path in $\overline\Om$ from
the maximizer to a boundary point, $u$ takes the value $v$. (The path lies in
$\overline\Om$, $u$ is continuous on it, and $u$ runs from $m$ to $0$.)
Since $\{u=v\}\subset\{u\ge v\}\subset\Om$, it lies in the open set $\Om$, and
$|\nabla u|>0$ on it because $v$ is a regular value.

\emph{Step 1: local analytic graphs (implicit function theorem).}
Fix $p\in\{u=v\}$. By Remark~\ref{rem:uniform-c2}, $u$ is real-analytic in a
neighbourhood of $p$, and $\nabla u(p)\ne0$. After relabelling coordinates we
may assume $\partial_{x_2}u(p)\ne0$. The real-analytic implicit function
theorem (\cite[Theorem~6.1.2]{KrantzParks2002}; for the $C^1$ statement
\cite[\S0.3 and Appendix]{GT1983}) yields an open neighbourhood
$U_p\ni p$, an interval $I_p$, and a real-analytic function
$g_p\colon I_p\to\R$ such that
\[
  \{u=v\}\cap U_p=\{(x_1,g_p(x_1)):x_1\in I_p\}.
\]
Thus near $p$ the level set is the graph of an analytic function, hence an
embedded analytic $1$-dimensional submanifold; the map $x_1\mapsto(x_1,g_p(x_1))$
is an analytic chart. As $p$ was arbitrary, $\{u=v\}$ is a real-analytic
embedded $1$-manifold without boundary in $\Om$. (No boundary: every point has
a chart modelled on an open interval, with no half-line charts, because the
graph extends to both sides of $p$ inside $U_p$.)

\emph{Step 2: finitely many components.}
By Step~1 each $p\in\{u=v\}$ has an open neighbourhood $U_p$ meeting $\{u=v\}$
in a single connected arc. The family $\{U_p\}_{p\in\{u=v\}}$ is an open cover
of the compact set $\{u=v\}$, so admits a finite subcover
$U_{p_1},\dots,U_{p_M}$. A connected component $\Gamma$ of $\{u=v\}$ is open
in $\{u=v\}$ (locally it is an arc, hence locally connected, so components are
open) and closed (components are always closed), thus compact; being covered by
finitely many of the $U_{p_j}$, and each component being a union of arcs that
overlap, there can be at most $M$ components, i.e. finitely many. Call them
$\Gamma_1,\dots,\Gamma_N$, $N\le M$.

\emph{Step 3: each component is an analytic Jordan curve.}
Each $\Gamma_i$ is a connected, compact, $1$-dimensional analytic manifold
without boundary. By the classification of connected $1$-manifolds
(\cite[Appendix, Classification of $1$-manifolds]{Milnor1997}), a connected
compact $1$-manifold without boundary is diffeomorphic to the circle $S^1$;
moreover the diffeomorphism can be taken analytic by reparametrising by analytic
arc length (the manifold being analytic). Hence each $\Gamma_i$ is the image of
an analytic embedding $\gamma_i\colon S^1\hookrightarrow\Om$, i.e. an analytic
Jordan curve. (It cannot be a line: a noncompact model is excluded by the
compactness of $\Gamma_i$ from Step~2.)

\emph{Step 4: $E_v$ open, bounded, with $\partial E_v=\{u=v\}$.}
$E_v=u^{-1}((v,\infty))\cap\Om$ is the preimage of an open set under the
continuous map $u\colon\Om\to\R$, hence open; it is bounded, being contained in
$\overline\Om$. For the boundary, write $A:=\{u>v\}$, $B:=\{u<v\}\cup(\R^2\setminus\overline\Om)$.
Both $A$ and the interior of $B$ are open and disjoint from $\{u=v\}$, so
$\partial E_v\subset\{u=v\}$: indeed any $x\notin\{u=v\}$ either has $u(x)>v$
(then $x\in A$ open, interior point of $E_v$) or has $u(x)<v$ or $x\notin\overline\Om$
(then a neighbourhood avoids $E_v$, so $x$ is exterior). Conversely let
$x\in\{u=v\}$. By Step~1, in $U_x$ the set $\{u=v\}$ is the graph of $g_x$, and
since $\partial_{x_2}u(x)\ne0$, the function $x_2\mapsto u(x_1,x_2)$ is strictly
monotone across the graph; hence both $\{u>v\}$ and $\{u<v\}$ meet every
neighbourhood of $x$. Therefore $x\in\partial E_v$. Combining,
$\partial E_v=\{u=v\}$.

\emph{Step 5: finite perimeter, perimeter value, and normal.}
By Step~4, $\partial E_v=\{u=v\}$ is a compact $C^1$ (indeed analytic)
hypersurface in $\R^2$ contained in $\Om$, with $E_v$ lying locally on one side
(the side where $u>v$). A bounded open set with $C^1$ boundary is a set of
finite perimeter, and its reduced boundary coincides with its topological
boundary, with measure-theoretic outer normal equal to the classical outer
normal and
\[
  P(E_v)=\HH^1(\partial E_v)=\HH^1(\{u=v\})=\sum_{i=1}^N\HH^1(\Gamma_i)
\]
(\cite[\S5.7.2, Theorem~1 and \S5.1]{EvansGariepy2015}; this is also the
content of (P4)). It remains to identify the outward normal. The classical
outward unit normal of $E_v=\{u>v\}$ at $x\in\partial E_v$ points in the
direction of \emph{decreasing} $u$, i.e. opposite to $\nabla u$, since $\nabla u$
points in the direction of steepest increase of $u$. Normalising,
\[
  \nu(x)=-\frac{\nabla u(x)}{|\nabla u(x)|},
\]
which is well-defined because $|\nabla u(x)|>0$ on the regular level set. This
proves (a)--(d).
\end{proof}

\subsection*{The flux identity}

\begin{lemma}[Flux identity]\label{lem:flux}
For every regular value $v\in(0,m)$, with $\mu(v):=|E_v|$,
\[
  \int_{\{u=v\}}|\nabla u|\,d\HH^1=|E_v|=\mu(v),
  \qquad
  \int_{\{u=v\}}\frac{1}{|\nabla u|}\,d\HH^1=-\mu'(v),
\]
the second identity holding for a.e.\ regular $v$ (at which $\mu$ is
differentiable).
\end{lemma}

\begin{proof}
\emph{First identity.}
Fix a regular value $v$. By Lemma~\ref{lem:smooth-level}, $\overline{E_v}\Subset\Om$
is compact with smooth (analytic) boundary $\partial E_v=\{u=v\}$, and by
Remark~\ref{rem:uniform-c2} $u\in C^\infty$ in an open neighbourhood of
$\overline{E_v}$. The vector field $X:=\nabla u$ is therefore $C^1$ up to
$\overline{E_v}$, and the divergence theorem
(\cite[\S4 and the Gauss--Green theorem]{EvansGariepy2015};
\cite[(2.10)--(2.11)]{GT1983}) applies on the bounded $C^1$
domain $E_v$:
\[
  \int_{E_v}\operatorname{div}(\nabla u)\,dx
  =\int_{\partial E_v}\nabla u\cdot\nu\,d\HH^1 .
\]
On the left, $\operatorname{div}(\nabla u)=\Delta u=-1$, so the left side equals
$-|E_v|$. On the right, $\nu=-\nabla u/|\nabla u|$ by
Lemma~\ref{lem:smooth-level}(d), hence
\[
  \nabla u\cdot\nu=\nabla u\cdot\Bigl(-\frac{\nabla u}{|\nabla u|}\Bigr)
  =-\frac{|\nabla u|^2}{|\nabla u|}=-|\nabla u|,
\]
that is $\partial_\nu u=-|\nabla u|$ on $\{u=v\}$. Therefore
\[
  -|E_v|=\int_{\{u=v\}}\bigl(-|\nabla u|\bigr)\,d\HH^1,
\]
and multiplying by $-1$ gives
$\int_{\{u=v\}}|\nabla u|\,d\HH^1=|E_v|=\mu(v)$, the first identity.

\emph{Second identity.}
We use the coarea formula (P4): for the locally Lipschitz function $u$ on the
open set $\{u>0\}$ and any Borel $g\ge0$,
\[
  \int_{\{u>0\}}g(x)\,|\nabla u(x)|\,dx
  =\int_0^{\infty}\Bigl(\int_{\{u=t\}}g\,d\HH^1\Bigr)dt
  \qquad(\cite[\S3.4.2, Theorem~5]{EvansGariepy2015}).
\]
Apply this with $g=\mathbf 1_{\{v<u\le m\}}\,|\nabla u|^{-1}$ on the set
$\{0<u\}$, noting $|\nabla u|^{-1}$ is finite $\HH^1$-a.e.\ on $\{u=t\}$ for
a.e.\ regular $t$ and that $g|\nabla u|=\mathbf 1_{\{v<u\le m\}}$. Equivalently,
choosing $g=\mathbf 1_{\{u>v\}}/|\nabla u|$ where defined, the left side is
$|\{u>v\}|=\mu(v)$ minus boundary contributions of measure zero; the cleanest
route is to differentiate the coarea representation of $\mu$. Indeed, taking
$g=\mathbf 1_{\{u>v\}}$ in the coarea formula written for the measure of
superlevel sets gives, for $0<v<m$,
\[
  \mu(v)=|\{u>v\}|=\int_{\{u>v\}}\!1\,dx
  =\int_{\{u>v\}}\frac{|\nabla u|}{|\nabla u|}\,dx
  =\int_v^{m}\Bigl(\int_{\{u=t\}}\frac{1}{|\nabla u|}\,d\HH^1\Bigr)dt,
\]
where in the last step we applied coarea with $g=\mathbf 1_{\{u>v\}}/|\nabla u|$
and used that $|\{\nabla u=0\}|=0$ (P3/Sard implies the critical \emph{values}
have measure zero; the critical \emph{set} $\{\nabla u=0\}$ has Lebesgue measure
zero because $u$ is real-analytic and nonconstant, so $\{\nabla u=0\}$ is a
proper analytic subvariety, hence $\HH^2$-null). Thus $\mu$ is, on $(0,m)$, the
integral from $v$ to $m$ of the function
\[
  v\mapsto \Phi(v):=\int_{\{u=v\}}\frac{1}{|\nabla u|}\,d\HH^1 ,
\]
which is finite for a.e.\ regular $v$. By the Lebesgue differentiation theorem
(fundamental theorem of calculus for the absolutely continuous function
$v\mapsto\int_v^m\Phi$), for a.e.\ $v$,
\[
  \mu'(v)=-\Phi(v)=-\int_{\{u=v\}}\frac{1}{|\nabla u|}\,d\HH^1 ,
\]
i.e. $\int_{\{u=v\}}|\nabla u|^{-1}\,d\HH^1=-\mu'(v)$, the second identity. The
minus sign reflects that $\mu$ is non-increasing (Lemma~\ref{lem:critvals}).
\end{proof}

\subsection*{Critical values and regularity of $\mu$}

\begin{lemma}[Finitely-thin critical set; regularity of $\mu$]\label{lem:critvals}
Let $\mathrm{Crit}:=\{v\in(0,m):v\text{ is a critical value of }u\}$. Then:
\begin{enumerate}
\item[(a)] $\mathrm{Crit}$ is a closed subset of $(0,m)$ of Lebesgue measure
  zero; consequently a.e.\ $v\in(0,m)$ is a regular value, and the conclusions of
  Lemmas~\ref{lem:smooth-level}--\ref{lem:flux} hold for a.e.\ $v$;
\item[(b)] for \emph{every} $v\in(0,m)$ the level set $\{u=v\}$ is Lebesgue-null:
  $|\{u=v\}|=\HH^2(\{u=v\})=0$;
\item[(c)] $\mu(v)=|E_v|$ is non-increasing on $(0,m)$ and continuous, with no
  jumps (no atoms).
\end{enumerate}
\end{lemma}

\begin{proof}
\emph{(a) Measure zero and closedness.}
The critical \emph{values} of $u$ have measure zero by Sard's theorem (P3),
applicable since $u\in C^2(\Om)$ (indeed $C^\infty$, by (P1)--(P2)). For
closedness: let $v_k\in\mathrm{Crit}$, $v_k\to v\in(0,m)$. Choose
$x_k\in\{u=v_k\}$ with $\nabla u(x_k)=0$. By
Lemma~\ref{lem:compact-containment}, all $x_k$ lie in the fixed compact set
$\{u\ge \tfrac{1}{2}\min_k v_k\}\Subset\Om$ (and $\min_k v_k>0$ since
$v_k\to v>0$); passing to a subsequence, $x_k\to x_*\in\Om$. By continuity of
$u$ and of $\nabla u$ (P2), $u(x_*)=v$ and $\nabla u(x_*)=0$, so $v$ is a
critical value, i.e. $v\in\mathrm{Crit}$. Hence $\mathrm{Crit}$ is closed in
$(0,m)$.

\emph{(b) Level sets are area-null.}
Two cases. If $v$ is a regular value, then by Lemma~\ref{lem:smooth-level}
$\{u=v\}$ is a finite union of analytic curves, each of which has zero
$2$-dimensional Lebesgue measure (a $1$-manifold in $\R^2$ is locally a graph,
hence $\HH^2$-null), so $|\{u=v\}|=0$. If $v$ is a critical value, write
\[
  \{u=v\}=\bigl(\{u=v\}\cap\{\nabla u\ne0\}\bigr)\cup\bigl(\{u=v\}\cap\{\nabla u=0\}\bigr).
\]
At points of the first set the implicit function theorem (as in
Lemma~\ref{lem:smooth-level}, Step~1) again represents $\{u=v\}$ locally as an
analytic graph, hence that part is $\HH^2$-null. The second set is contained in
$\{\nabla u=0\}$, which is Lebesgue-null because $u$ is real-analytic and
nonconstant on the connected components of $\Om$ (so $\nabla u$ does not vanish
identically): the zero set of a nonzero real-analytic function has Lebesgue
measure zero (\cite[Lemma in \S0; or the analytic zero-set theorem]{KrantzParks2002}).
Covering $\{u=v\}$ by countably many balls on which $u$ is analytic and applying
this in each, we get $|\{u=v\}|=0$ for every $v$.

\emph{(c) Monotonicity and continuity of $\mu$.}
Monotonicity: if $v_1<v_2$ then $E_{v_2}=\{u>v_2\}\subset\{u>v_1\}=E_{v_1}$, so
$\mu(v_2)=|E_{v_2}|\le|E_{v_1}|=\mu(v_1)$; thus $\mu$ is non-increasing.
Continuity: for $v_1<v_2$ in $(0,m)$,
\[
  \mu(v_1)-\mu(v_2)=|E_{v_1}\setminus E_{v_2}|=|\{v_1<u\le v_2\}|.
\]
As $v_2\downarrow v_1$, the sets $\{v_1<u\le v_2\}$ decrease to
$\{u=v_1\}\cup\varnothing=\{u=v_1\}$ (more precisely
$\bigcap_{v_2>v_1}\{v_1<u\le v_2\}=\{u=v_1\}$), which has measure $0$ by part
(b); by dominated convergence (all sets contained in the finite-measure set
$\Om$) the difference tends to $0$, giving right-continuity. Similarly, as
$v_1\uparrow v_2$ the sets $\{v_1<u<v_2\}$ decrease to $\varnothing$, while
$\{v_1<u\le v_2\}\downarrow\{u=v_2\}$, again of measure $0$, giving
left-continuity. Hence $\mu$ is continuous on $(0,m)$ with no atoms.
\end{proof}

\begin{remark}\label{rem:ae-suffices}
The combination of Lemmas~\ref{lem:smooth-level}--\ref{lem:critvals} is exactly
what later sections require: away from the measure-zero set $\mathrm{Crit}$,
the superlevel sets $E_v$ are smooth finite-perimeter domains for which the
flux identity and the coarea representation of $\mu$ hold pointwise, while the
exceptional set $\mathrm{Crit}$ may be discarded in any integration in $v$
(e.g. when integrating the flux identity against $dv$). When $u$ is moreover
real-analytic, $\mathrm{Crit}$ is in fact finite, but we shall only use that it
is closed and null.
\end{remark}
